\documentclass[12pt]{article}
\usepackage[spanish]{babel} 
\usepackage{packages/mygeneral}
\usepackage{packages/mycommands}
\usepackage{packages/mylayout} 
\institution{
  name=UAM
} 
\pagelayout{
 rightheader=Física Atómica y Molecular,
 logo=false % set false/true to hide/show the lowo in the left header
} 
\titlesetup{
  theme=project,
  title=Método optimizado de obtención de términos L-S y j-j para electrones degenerados,
  author={Joan A. Mercado Tandazo $\quad$ Juan Manuel Sánchez Arrúa},
  projectname=Física Atómica y Molecular
} 
\mathtoolsset{showonlyrefs=true} % for showing number only in labeled equations
\allowdisplaybreaks[2] %in case you want your equations to be broken across pages
\let\vec\mathbf % display vectors in bold instead of with an arrow
\begin{document}
\maketitle
\tableofcontents
\resetnumpagetitle
\pagestyle{academic} 
% !TEX root = main.tex 
\section{Introducción} 
El método que se presenta a continuación es una mejora del tradicional que se presenta en el curso de Física Atómica y Molecular, el cual suele estar basado en tablas y contabilidad de posibilidades según estas se le ocurran al estudiante. El método que presentamos reduce el esfuerzo de contabilidad gracias a la combinatoria y simetrías de los términos, sistematizando el proceso y siendo fácilmente extensible a un número arbitrario de electrones.\par
El procedimiento consiste en centrarse en el número cuántico $M_{L}$ (o $M_{J}$ para el método j-j) e ir obteniendo los términos de acuerdo al número de combinaciones posibles de los electrones para dar ese número cuántico. Conviene empezar desde el máximo valor posible de $M_{L}$ (o $M_{J}$) e ir reduciendo su valor hasta llegar a $M_{L}=0$ (o $M_{J}=0$) (hasta ese valor será suficiente gracias a las simetrías). Esta reducción se puede realizar sistemáticamente gracias a la aplicación del operador $L_{-}$ (o $J_{-}$) como veremos en los ejemplos. Según las distintas combinaciones que surjan, y teniendo en cuenta el spin para el caso de L-S, obtendremos los términos correspondientes, que deberemos ir anotando a medida que los vayamos encontrando, puesto que la siguiente reducción los necesitará en la contabilidad. 
 
% !TEX root = main.tex
\section{Términos L-S}\label{sec: terminos L-S}
Supongamos que tenemos $N$ electrones en un mismo subnivel $l$, y queremos obtener los términos L-S correspondientes. Es decir, partimos de una configuración electrónica de la forma $l^{N}$ donde asumimos que $N \leq (1/2)(2l + 1)\times (2s+1)$, pues por la simetría ``electrón-hueco'' el caso de $N > (1/2)(2l + 1)\times (2s+1)$ se puede obtener a partir del caso de $N' = (2l + 1)\times (2s+1) - N$ electrones. Por teoría del momento angular, sabemos que el número cuántico total de momento orbital $L$ puede tomar valores entre $L= N \times l$ y $ L =0$ ó $L=l$ según $N$ sea par o impar respectivamente. Esto nos da un primer acercamiento de los valores de $L$ que podemos esperar en los términos L-S. \par
Una parte fundamental del método es estudiar el spin. Análogamente al momento angular, el spin total tomará valores entre $S = N/2$ y $S = 0$ ó $S = 1/2$ según $N$ sea par o impar. Como las posibilidades de la tercera componente del spin individual son menores a las del momento orbital (individual) entonces le asignamos a este un rol secundario. De esta forma tenemos el siguiente panorama: para cada combinación $(M_{L},M_{S})$ tenemos un conjunto de posibles <<tuplas>>\footnote{Nótese que en general no hablamos de una correspondencia entre estados $\ket{M_{L},M_{S}}$ con estados $\ket{m_{1}, \sigma_{1};m_{2},\sigma_{2};\dots ; m_{N}, \sigma_{N}}$, pues en general $\ket{M_{L}, M_{S}}$ será una combinación lineal de estos últimos.}  de los $e^{-}$ individualmente $(m_{1}^{\sigma_{1}}, m_{2}^{\sigma_{2}},\dots m_{N}^{\sigma_{N}})$ donde $m_{i}$ es la tercera componente del momento orbital del electrón $i$ y $\sigma_{i}$ es su tercera componente del spin.  Además, por teoría del momento angular deben cumplir: 
\begin{align}
  \sum_{i=1}^{N} &m_{i} = M_{L} & \sum_{i=1}^{N} &\sigma_{i} = M_{S} \label{eq: suma de 3as componentes} 
\end{align}
Cada uno de estos $m_{i}^{\sigma_{i}}$ son distintos para cumplir el principio de exclusión de Pauli. Existen $2\times(2l + 1))$ posibilidades de $m_{i}^{\sigma_{i}}$ para cada electrón, y por lo tanto, el número de formas posibles de tomar $N$ electrones distintos de entre estas $2\times(2l + 1))$ posibilidades es:
\begin{align}
\# \text{combinaciones} = \binom{2 \times (2l+1)}{N}\label{eq: numero de combinaciones final}
\end{align}
Lo cual se puede comprobar que coincide con el resultado del método tradicional, pues de hecho, es una explicación de dicho resultado. Esto también nos quiere decir que la degeneración acumulada de los términos L-S que obtendremos ha de coincidir con este número de combinaciones, lo cual será una prueba de verificación del método. \par 

A continuación centrémonos en $M_{L}$ y busquemos el número más alto posible que puede tener. Por \eqref{eq: suma de 3as componentes} tendremos que buscar los máximos valores posibles de $m_{i}$ respetando, no obstante, el principio de exclusión. Para ello introducimos la siguiente regla general que nos será útil para el resto del método:\par 
\begin{center} 
\fbox{
  \begin{minipage}[t]{0.9\textwidth}
  Denominando a la expresión $M_{L} = m_{1} + m_{2} + \dots + m_{N}$ como una \textbf{partición} de $M_{L}$, entonces, en ella no podemos repetir el mismo $m$ 2 veces.
  \end{minipage}
}
\end{center}
\vspace*{0.2cm} 
La partición obtenida para el máximo valor de $M_{L}$ se corresponderá con el máximo valor de $L$. \par 
Así, nos queda averiguar el spin. Para ello, deberes sumar las $\sigma$'s de las cuales, por supuesto, no hemos llevado su contabilidad, porque no es necesario. Para la partición de $M_{L}$ más alto, es sencillo averiguar $M_{S}$ porque tendremos muchas repeticiones de $m$'s lo cual significa que las terceras componentes de spin se anulan entre sí y solo tendremos que fijarnos en la $m$ no repetida para averiguar $M_{S}$. Es decir, renombrando los valores que toma $\sigma$ como $+1/2 \to  +$ $-1/2 \to  -$ tendremos una forma tal que: 
\begin{align*}
  M_{L} &= \underbrace{a^{+} + a^{-} + b^{+} + b^{-} + \dots (+ m^{+})}_\text{$N$ sumandos}\; , \quad \text{ donde } a=\max(m_{i})>b>\dots > m 
\end{align*}
Nótese que asignamos los $+$ antes que los $-$, lo cual es una convención, pero no nos hace falta considerar el signo opuesto porque el resultado será simétrico y será irrelevante para obtener $S$. \par
Para $N$ par, no existirá esa $m$ no repetida. Entonces $M_{S} = 0$, lo cual inducirá a que $S=0$ para ese término de máxima $L$. Y para el caso de $N$ impar, existirá ese término de $m$ no repetida y entonces $M_{S} = 1/2$, lo cual inducirá a que $S=1/2$ para ese término de máxima $L$. De ahí podremos extraer el término L-S correspondiente con su degeneración, que será importante para la contabilidad de los siguientes términos. \par
Aunque en general se puede deducir para cada ejemplo, como nota de interés, el lector puede comprobar que en general el valor máximo de $L$, dado  $N$ y $l$ se puede obtener teniendo en cuenta que un valor dado de $m_{l}$ puede como máximo repetirse dos veces. De este modo: 
\begin{align}
   L_\text{max} =\underbrace{2l}_{1} + \underbrace{2(l-1)}_{2}+\ldots+\underbrace{2\left(l-\text{floor}(\frac{N}{2})-1 \right)}_{\text{floor}\left(\frac{N}{2}\right) }+\left[l-\text{floor}\left(\frac{N}{2}\right)-1-1\right]\times \text{mod}(N,2)\\
   L_\text{max} = 2l\times \text{floor}\left(\frac{N}{2}\right)-2\frac{\left(\text{floor}\left(\frac{N}{2}\right)+1\right)\left(\text{floor}\left(\frac{N}{2}\right)+2\right)}{2}+\left[l-\text{floor}\left(\frac{N}{2}\right)-2\right]\times \text{mod}(N,2)   
 \end{align}

Para obtener los siguientes términos, formalmente lo siguiente es aplicar el operador $L_{-}$ para reducir en 1 el valor de $M_{L}$. Si bien recordamos, dicho operador viene dado por: 
\begin{align*}
L_{-} = \sum_{i=1}^{N} l_{-}^{(i)}
\end{align*}
Y si aplicamos $L_{-}$ a un estado $\ket{m_{1},\sigma_{1}; m_{2},\sigma_{2}; \dots m_{N},\sigma_{N}}$ con $M_{L} = \sum_{i=1}^{N} m_{i}$ definido (aunque no tenga un $L$ definido), entonces el resultado será un estado con $M_{L} - 1$. Pero además: 
\begin{align*}
L_{-} \ket{m_{1},\sigma_{1}; m_{2},\sigma_{2}; \dots m_{N},\sigma_{N}} = \sum_{i=0}^{N} l_{-}^{(i)} \ket{m_{1},\sigma_{1}; m_{2},\sigma_{2}; \dots m_{N},\sigma_{N}} = \sum_{i=0}^{N} a_{i} \ket{\dots m_{i} - 1, \sigma_{i} \dots}
\end{align*}
Donde $a_{i}$ por construcción será tal que cumpla el principio de exclusión de Pauli, es decir, $a_{i} = 0$ si $m_{i} - 1$ ya está repetido. \par 
En el lenguaje simple de particiones esto se traduce en que a cada sumando de la partición anterior podemos reducirlo en una cifra para obtener una partición de $M_{L}-1$, siempre y cuando esa reducción no implique repetir más de 2 veces un sumando. Esto dará lugar a varias particiones y tras varias reiteraciones, es posible que a la hora de seguir aplicando $L_{-}$ se obtenga una partición repetida. En ese caso, no la contamos 2 veces, simplemente la ignoramos. Así podremos generar todas las particiones de $M_{L}-1$ a partir de las particiones de $M_{L}$, y sistemáticamente llegar a $M_{L} = 0$. \par 
Para cada uno de los $M_{L}$ tendremos que buscar los posibles $M_{S}$ correspondientes. En ese caso tendremos que generalizar lo que vimos para el caso de $M_{L}$ máximo. En estas particiones tendremos sumandos que se repiten y otros que no. De nuevo, para obtener $M_{S}$ solo nos interesa los términos no repetidos, puesto que los repetidos anulan la tercera de componente de spin entre sí. Asumamos que existen $k$ sumandos no repetidos y sabemos que el número de $M_{S}>0$ distintos será $\mathrm{floor}\qty(k/2) + 1$ porque cada sumando no repetido puede ser $+$ o $-$. Entonces, para una cierta partición $j$, nos interesa conocer el número de combinaciones distintas de $(+)$ y $(-)$ que podemos asignar a los $k$ sumandos no repetidos para un cierto $M_{S}>0$ dado, pues será un indicativo de a qué término $L-S$ corresponde. Para conocer ese número de combinaciones o degeneración, $G^{(j)}_{k}(M_{S})$, simplemente aplicamos combinatoria: 
\begin{align}
  G^{(j)}_{k}(M_{S}) &= \binom{k}{k/2 - M_{S}} \text{ con } k/2 \geq M_{S} > 0 \label{degeneracion M_S}\\ 
  \text{ Por completitud: } G^{(j)}_{k}(M_{S}) &= 0 \text{ para } M_{S} > k/2
\end{align}
Que se puede entender de la siguiente forma: podemos centrarnos en los $(-)$ que debe haber para dar un cierto $M_{S}>0$, entonces, podemos pensar en este resultado como las distintas formas que hay de asignar los $(-)$ a los $k$ sumandos no repetidos. Debemos tener en cuenta que $0! = 1$ es decir que $G_{k}(k/2) = 1$.\par 
Así, para cada partición de $M_{L}$ podremos tener una o varias posibilidades de $M_{S}$. Y esto es esperable, puesto que no solo un cierto $M_{L}$ puede tener varios $L$ asociados, los cuales pueden tener un $S$ distinto; sino que también podemos tener para un mismo $L$ con distintos $S$ asociados. El conteo de particiones de $M_{L}$ con sus correspondientes degeneraciones $G_{k}(M_{S})$ nos ha de dar como mínimo la degeneración total de los términos L-S obtenidos en las particiones de $M_{L}$ anterior y posiblemente, un excedente que se corresponden con términos L-S nuevos de $L$ menor.\par 
Debe tenerse en cuenta que se pueden encontrar términos L-S nuevos repetidos. No es ningún error, es simplemente que hay varias configuraciones de los electrones que pueden dar el mismo término L-S. Estos términos los marcaremos como $K \times L^{S}$ donde $K$ es el número de veces que se ha encontrado ese término L-S. \par 
Merece la pena introducir el concepto paralelo de \textit{número de duplas iguales} $\mathcal{D}$, a las que nos referiremos simplemente como \textit{duplas}. Estas no son más que el número de parejas iguales en una cierta partición de $M_{L}$. Se encuentran relacionadas con el número de sumandos no repetidos $k$ de la siguiente forma: 
\begin{align*}
N = 2 \mathcal{D} + k 
\end{align*}
Donde $N$ es el número de electrones. En los ejemplos que comentemos en los puntos siguientes emplearemos las duplas para establecer \textbf{reglas} que faciliten la contabilidad y obtención de los términos L-S.\par
Finalmente, el método se da por terminado cuando llegamos a $M_{L} = 0$. En ese caso, la degeneración total del número de términos L-S encontrados debe coincidir con el número de combinaciones \eqref{eq: numero de combinaciones final}, lo cual es una prueba de que el método se ha aplicado correctamente. \par


 
\subsection{Ejemplo: 2 \texorpdfstring{$e^{-}$}{e^-}} 
% !TEX root = main.tex

\renewcommand{\arraystretch}{1.5}
Consideramos el caso de 2 $e^-$ con  $l=1$. Las posibles 
configuraciones electrónicas serán aquellas que 
obedezcan el Principio de Exclusión de Pauli. Para 
encontrarlas, vamos a considerar las posibles 
configuraciones de espín  en función del 
número de duplas $(m_{l}^{(1)},m_{l}^{(2)})$ con $m_{l}^{(1)}=m_{l}^{(2)}$.
En este caso, al tener solo 2 $e^-$, el número máximo 
de duplas iguales es 1, pero en general dependerá del valor de  $l$. 
%NOTE:(Se puede buscar cuál es el nro. máximo de duplas que se pueden formar en función de $l$ y  $N$).
Una vez determinamos el número máximo de duplas 
permitido $\mathcal{D}_{M}$, es fácil ver que van a 
existir configuraciones permitidas con valor de
duplas $\mathcal{D}\in [0,\mathcal{D}_{M}]$. A cada valor de dupla $\mathcal{D}$ 
le es asociado un conjunto de posibles configuraciones de 
espín.\par 
Para el caso particular de $2$  $e^-$ con  $l=1$, 
tenemos  $\mathcal{D}_{M}=1$ y las \textbf{<<reglas>>} o 
posibles configuraciones de espín para 
cada $\mathcal{D}$ son:\par 
\underline{\textbf{Reglas}}
\begin{align*}
  \mathcal{D}=1&\longrightarrow \begin{cases}
    1\times M_{s}=0:\quad (+,-)
  \end{cases}\\
    \mathcal{D}=0&\longrightarrow \begin{cases}
      1\times M_{s}=1:\quad (+,+)\\
      2\times M_{s}=0:\quad (+,-)
    \end{cases}
\end{align*}
Es quizás importante remarcar, en el caso de $\mathcal{D}=1$, 
que las configuraciones $(+,-)$ y $(-,+)$ no son distintas 
y sólo se cuenta una de ellas. Esto es porque en este
caso  $m_{l}^{(1)}=m_{l}^{(2)}$, y como el resto de 
números cuánticos también son iguales, los 
$e^-$ son equivalentes, y el estado 
de esta configuración debe ser la combinación antisimétrica de las configuraciones mencionadas 
 \begin{align}
  \ket{\psi} = \frac{1}{\sqrt{2}}\left(\ket{n,n;l,l;m_{l},m_{l};+,-}-\ket{n,n;l,l;m_{l},m_{l};-,+} \right) 
\end{align}
es decir ambas configuraciones refieren al mismo estado. 
Estamos empleando la siguiente notación para 
especificar una configuración 
\begin{align*}
  \ket{\psi} = \ket{n_{1},n_{2};l_{1},l_{2};m_{l}^{(1)},m_{l}^{(2)};m_{s}^{(1)},m_{s}^{(2)}}
\end{align*}
Mientras que en el caso $\mathcal{D}=0$, al tener $m_{l}$ distintos, 
los dos electrones son inequivalentes, y por lo tanto 
cuando $M_{S}=0$ tenemos dos posibles estados diferentes 
\begin{align*}
  \ket{\psi_{1}} &= \frac{1}{\sqrt{2}}\left(\ket{n,n;l,l;m_{l},m_{l}';+,-}-\ket{n,n;l,l;m_{l}',m_{l};-,+}\right)\\
 \ket{\psi_{2}} &= \frac{1}{\sqrt{2}}\left(\ket{n,n;l,l;m_{l},m_{l}';-,+}-\ket{n,n;l,l;m_{l}',m_{l};+,-}\right)\\
  \ket{\psi_{1}}&\ne \ket{\psi_{2}}
\end{align*}
Establecidas las reglas, procedemos a obtener los términos según el valor de $M_{L}$:
\begin{align*}
 \begin{array}{c | c | c | c}
  M_{L} & \text{ Particiones } & \text{ Duplas } & \text{ Términos }\\
  \hline 
  M_{L} = 2 & 2 = 1 + 1 & \mathcal{D} = 1,\; k = 0 & 1 \times (M_{S} = 0) \implies \ce{^{1}D}(5)\\
  \hline 
  M_{L} = 1 & 1 = 1 + 0 & \mathcal{D} = 0,\; k = 2 & \begin{cases}
    1 \times (M_{S} = 1) \implies \ce{^{3}P}(9)\\
    2 \times (M_{S} = 0) \implies \cancel{\ce{^{1}D}(5)},\ \cancel{\ce{^{3}P}(9)}
    \end{cases}\\
    \hline 
  M_{L} = 0 & 0 = 1 - 1 & \mathcal{D} = 0,\; k = 2 & \begin{cases}
    1 \times (M_{S} = 1) \implies \cancel{\ce{^{3}P}(9)}\\
    2 \times (M_{S} = 0) \implies \cancel{\ce{^{1}D}(5)},\ \cancel{\ce{^{3}P}(9)}
    \end{cases}\\
  & 0 = 0 + 0 & \mathcal{D} = 1,\; k = 0 & 1 \times (M_{S} = 0) \implies \ce{^{1}S}(1)\\
  \hline
  \end{array}
  \end{align*}
Una de las ventajas de este método es que, a la hora 
de determinar las posibles combinaciones de $m_{l}$ que 
suman $M_{L}$ basta con partir del $M_{L}^{(\text{max})}$ máximo que 
se obtiene trivalmente de maximizar todos los $m_{l}$ 
(respetando el principio de Pauli), y luego para 
$M_{L}^{(\text{max})}-1$ las combinaciones posibles se obtienen 
de restar $1$ a los diferentes $m_{l}$ que aparecen 
en la configuración de $M_{L}^{(\text{max})}$.\par 
Por otro lado, al disminuir $M_{L}$ y contar los 
estados, tenemos que considerar que para una 
configuración L-S dada, también van a 
existir instancias con  $M_{L}$ y $M_{S}$ menores, 
por lo que en cada $M_{L}$ debemos de asegurarnos 
de no perder estos estados ni contarlos como 
configuraciones nuevas o distintas. Específicamente a tener en cuenta en la tabla anterior es que, en cada $M_{L}$ hemos contabilizado y posteriormente tachado los términos que ya han surgido en $M_{L}$ mayores, y que, por lo tanto, no son nuevos.\par

Para $M_{L}<0$ las configuraciones serán 
idénticas, simplemente cambiando el signo de las $m_{l}$. 
Entonces, los términos espectrales que obtenemos son:
\begin{align*}
  \Aboxed{\text{T.E.:}\quad &\ce{^{1}D}(5), \quad \ce{^{3}P}(9),\quad \ce{^{1}S}(1)}\\ 
  \text{ Degeneración total: }&g_{T} = 5 + 9 + 1 = 15 = \binom{6}{2}\quad \boxed{\checkmark}
\end{align*}
 

\subsection{Ejemplo: 3 \texorpdfstring{$e^{-}$}{e^-}} 
% !TEX root = main.tex
\renewcommand{\arraystretch}{1.5}
Consideremos ahora el caso de 3 $e^{-}$ degenerados con $l=2$. La correspondiente configuración electrónica es $\ce{d^{3}}$. Ahora tendremos que considerar triadas del tipo $\qty(m_{1}^{\sigma_{1}}, m_{2}^{\sigma_{2}},m_{3}^{\sigma_{3}})$. En cuanto al número de duplas que se pueden formar, es fácil observar que se pueden formar $\mathcal{D} \in  \qty{0,1}$ duplas, que se corresponden con $k \in \qty{3,1}$ sumandos no repetidos. Entonces, mediante la degeneración $G_{k}$ por partición podemos establecer las siguientes reglas:\par
\underline{\bfseries Reglas}:
\begin{align*}
  \mathcal{D}&=1 \iff k=1 \implies 1 \times \qty(M_{S} = 1/2)\\ 
  \mathcal{D}&=0 \iff k=3 \implies 
  \begin{dcases}
    1 \times \qty(M_{S} = 3/2): & \qty(+, +, +)\\
    3 \times \qty(M_{S} = 1/2): & \text{permutaciones de } \qty(+, +, -)
  \end{dcases} 
\end{align*}
Con estas reglas procedemos a desarrollar el proceso reiterativo (veáse el cálculo adjunto), primero buscando el $M_{L}$ máximo que resulta en $5$. Como el número de electrones es impar, $M_{S} = 1/2$ para este término, resultado en $\ce{^{2}H}(22)$. A continuación, aplicamos $L_{-}$, i.e., reducimos en una cifra uno de los sumandos de la anterior partición respetando el principio de exclusión. De ahí obtenemos 2 particiones, ambas con $\mathcal{D} = 1$, $k=1$, lo cual nos lleva inmediatamente a $M_{S} = 1/2$  en ambas. Deducimos entonces la presencia de dos términos, uno con $L = 5$ (el anterior término L-S) y otro con $L = 4$, resultando en $\ce{^{2}G}(18)$. Nótese como los términos ya obtenidos, los contabilizamos y tachamos para indicar que no son nuevos. A continuación, volvemos a reiterar el método. Para $M_{L} = 3$ nos encontramos con una partición con $\mathcal{D} = 0$, $k=3$, donde justamente las ventajas de este método salen a brillar. Con las reglas que hemos establecido gracias a la combinatoria presentada, deducimos en esta partición tenemos un término con $M_{S} = 3/2$ y tres términos con $M_{S} = 1/2$. El término con $M_{S} = 3/2$ es un término L-S nuevo, el cual corresponde a $\ce{^{4}F}(28)$. En cuanto a los términos con $M_{S} = 1/2$, podemos ver que tenemos contabilizados los dos anteriores, pero además, el nuevo que hemos obtenido con $S = 3/2$ también tiene que ser contado aquí, pues su tercera componente puede tomar valores tanto de $M_{S= 3/2}$ como de $M_{S} = 1/2$. Este es un detalle general importante a tener en cuenta en todo el método y que garantiza que la contabilidad no se desajuste. Nótese que también encontramos un término $\ce{^{2}F}(14)$.\par 
En las reiteraciones siguientes, hemos omitido escribir todos los términos ya presentes, centrándonos solo en los nuevos. Simplemente aplicando la regla a cada una de las particiones y sumando todos los casos de $M_{S}=3/2$ y $M_{S} = 1/2$ que se vayan obteniendo, podemos inducir, si surgen nuevos términos de menor $L$ o no. Por ejemplo, para $M_{L} = 3$ obtuvimos $1 \times (M_{S} = 3/2)$ y $4 \times (M_{S} = 1/2)$, mientras que para $M_{L} = 2$  se ha obtenido $1 \times (M_{S} = 3/2)$ y $6 \times (M_{S} = 1/2)$. Esto nos lleva a deducir que para $M_{L} = 2$ aparecen dos términos nuevos con $S = 1/2$, es decir, $2 \times \ce{^{2}D}(10)$. Con la misma metodología, llegamos a que para $M_{L} =1$, que tiene $2 \times (M_{S} = 3/2)$  y $8 \times (M_{S} = 1/2)$, debe haber un término con $S = 3/2$, $\ce{^{4}P}(12)$, y un término nuevo con $S = 1/2$, i.e. $\ce{^{2}P}(6)$, (nótese cómo de nuevo que el término con $S = 3/2$ también se contabiliza en el caso de $M_{S} = 1/2$). Finalmente, para $M_{L} = 0$, no aparecen términos nuevos. 
\begin{align*}
 \begin{array}{c | c | c | c}
   M_{L} & \text{ Particiones } & \text{ Duplas } & \text{ Términos }\\
   \hline
   M_{L}= 5  & 5 = 2 + 2 + 1  &  \mathcal{D} = 1,\; k = 1 & 1 \times (M_{S} = 1/2) \implies \ce{^{2}H}(22)\\
   \hline
   M_{L} = 4 & 4 = 2 + 2 + 0 & \mathcal{D} = 1,\; k = 1 & 1 \times (M_{S} = 1/2) \implies \cancel{\ce{^{2}H}(22)}\\ 
             & 4 = 2 + 1 + 1 & \mathcal{D} = 1,\; k = 1 & 1 \times (M_{S} = 1/2) \implies \ce{^{2}G}(18)\\
  \hline
   M_{L} = 3 & 3 = 2 + 1 + 0 & \mathcal{D} = 0, \; k = 3 & \begin{cases}
   1 \times (M_{S} = 3/2) &\implies \ce{^{4}F}(28)\\
   3 \times (M_{S} = 1/2) & \implies \cancel{\ce{^{2}H}(22)}, \cancel{\ce{^{2}G}(18)}, \cancel{\ce{^{4}F}(28)}
   \end{cases}\\ 
             & 3 = 2 + 2 -1 & \mathcal{D} = 1, \; k = 1 & 1 \times (M_{S} = 1/2) \implies \ce{^{2}F}(14)\\
  \hline
     M_{L} = 2 & 
     \begin{aligned}
       2 &= 1 + 1 + 0\\ 
         &= 2 + 0 + 0\\ 
         &= 2 + 1 - 1\\ 
         &= 2 + 2 - 2
     \end{aligned} 
     & \begin{aligned}
       \mathcal{D} &= 1,\;  k = 1\\ 
       \mathcal{D} &= 1,\;  k = 1\\ 
       \mathcal{D} &= 0,\;  k = 3\\ 
       \mathcal{D} &= 1,\;  k = 1\\ 
     \end{aligned} & \begin{dcases}
      1 \times (M_{S} = 3/2) \\
      6\times (M_{S} = 1/2) 
     \end{dcases}  \implies 2 \times \ce{^{2}D}(10) \\ 
     \hline
     M_{L} = 1 & 
     \begin{aligned}
     1 &= 1 + 0 + 0\\ 
       &= 2 + 1 - 2\\ 
       &= 1 + 1 - 1\\ 
       &= 2 + 0 - 1
   \end{aligned} & 
   \begin{aligned}
      \mathcal{D} &= 1,\; k = 1\\ 
      \mathcal{D} &= 0,\; k = 3\\ 
      \mathcal{D} &= 1,\; k = 1\\ 
      \mathcal{D} &= 0,\; k = 3\\ 
    \end{aligned} & 
    \begin{cases}
    2 \times (M_{S} = 3/2)\\
    8 \times (M_{S} = 1/2) 
    \end{cases} \implies  \ce{^{4}P}(12), \ce{^{2}P}(6)\\
 \hline 
     M_{L} = 0 & 
     \begin{aligned}
     0 &= 1 + 1 - 2\\ 
       &= 2 + 0 - 2\\ 
       &= 1 + 0 - 1\\ 
       &= 2 -1 - 1
   \end{aligned} & 
   \begin{aligned}
      \mathcal{D} &= 1,\; k = 1\\ 
      \mathcal{D} &= 0,\; k = 3\\ 
      \mathcal{D} &= 0,\; k = 3\\ 
      \mathcal{D} &= 1,\; k = 1\\ 
    \end{aligned} & 
    \begin{cases}
    2 \times (M_{S} = 3/2)\\
    8 \times (M_{S} = 1/2) 
    \end{cases} \implies  \text{ Ningún término L-S nuevo }\\
 \end{array}
\end{align*}
La prueba de degeneración acumulada $g_{T}$ nos garantiza que los términos obtenidos son correctos: 
\begin{align*}
  g_{T} &= 22 + 18 + 28 + 14 + 2 \times 10 + 12 + 6 = 210 = \binom{2 \times 5}{3} & &\boxed{\checkmark}
\end{align*}
Por tanto, los términos espectrales (T.E.) son: 
\begin{align*}
  \boxed{\text{ T.E.:} \quad \ce{^{4}F}(28), \quad \ce{^{2}H}(22), \quad \ce{^{2}G}(18), \quad \ce{^{2}F}(14), \quad 2 \times \ce{^{2}D}(10), \quad \ce{^{4}P}(12), \quad \ce{^{2}P}(6)}
\end{align*}


 
\subsection{Ejemplo: 4 \texorpdfstring{$e^{-}$}{e^-}}
Consideramos por último, el caso de 4 $e^-$ equivalentes con $l=2$. La correspondiente configuración electrónica es $d^{4}$. En este caso $\mathcal{D}\in [0,2]$ y las reglas son \\
\underline{\textbf{Reglas}}:
\begin{align}
  &\mathcal{D}=2 \longrightarrow 1\times (M_{S}=0):\ (+-+-)\\
  &\mathcal{D}=1 \longrightarrow \begin{cases}
    1\times (M_{S}=1):\ (+-++)\\
    2\times (M_{S}=0):\ (+-+-)
  \end{cases}\\
  &\mathcal{D}=0 \longrightarrow \begin{cases}
    1\times (M_{S}) = 2:\ (++++)\\
    4\times (M_{S}=1):\ (+++-)\\
    6\times (M_{S}=0):\ (++--)
  \end{cases}
 \end{align}
Entonces los términos espectrales son 


\begin{align*}
 \begin{array}{c | c | c | c}
   M_{L} & \text{ Particiones } & \text{ Duplas } & \text{ Términos }\\
   \hline
   M_{L}= 6  & 6 = 2 + 2 + 1 + 1  &  \mathcal{D} = 2 & 1 \times (M_{S} = 0) \implies \ce{^{1}I}(13)\\
   \hline
   M_{L} = 5 & 5 = 2 + 2 + 1 + 0 & \mathcal{D} = 1 & \begin{dcases} 1 \times (M_{S} = 1) \implies \ce{^{3}H}(33)\\
     2\times (M_{S}=0)\implies \cancel{\ce{^{1}I}(13)},\ \cancel{\ce{^{3}H(33)}} 
   \end{dcases}\\
  \hline
     M_{L} = 4 & 
     \begin{aligned}
       4 &= 2 + 2  + 0 + 0\\ 
         &= 2 + 1 + 1 + 0\\
         &= 2 + 2 + 1 - 1
     \end{aligned} 
     & \begin{aligned}
       \mathcal{D} &= 2\\ 
       \mathcal{D} &= 1\\
       \mathcal{D} &=1
     \end{aligned} & \begin{dcases}
      2 \times (M_{S} = 1) \\
      5\times (M_{S} = 0) 
   \end{dcases}  \implies \ce{^{3}G(27)},\ 2\times\ce{^{1}G(9)} \\ 
     \hline
     M_{L} = 3 & 
     \begin{aligned}
     3 &= 2 + 1 + 0 + 0\\ 
       &= 2 + 1 + 1 - 1\\
       &= 2 + 2 - 1 + 0\\
       &= 2 + 2 + 1 -2 
   \end{aligned} & 
   \begin{aligned}
      \mathcal{D} &= 1\\ 
      \mathcal{D} &= 1\\
      \mathcal{D}&= 1\\
      \mathcal{D}&=1 
   \end{aligned} & 
    \begin{cases}
    4 \times (M_{S} = 1)\\
    8 \times (M_{S} = 0) 
  \end{cases} \implies  2\times\ce{^{3}F}(21),\ \ce{^{1}F}(7)\\
 \hline
   M_{L} = 2 & 
     \begin{aligned}
       2 &= 1 + 1 + 0 + 0\\ 
         &= 2 + 1 - 1 + 0\\
         &= 2 + 1 + 1 - 2\\
         &= 2 + 2 + 0 - 2\\
         &= 2 + 2 - 1 - 1
     \end{aligned} 
     & \begin{aligned}
       \mathcal{D} &= 2\\ 
       \mathcal{D} &= 0\\
       \mathcal{D} &= 1\\
       \mathcal{D} &= 1\\
       \mathcal{D} &= 2
     \end{aligned} & \begin{dcases}
      1 \times (M_{S} = 2) \\
      6\times(M_{S}=1)\\ 
      12\times (M_{S} = 0) 
   \end{dcases}  \implies \ce{^{5}D(25)}, \ce{^{3}D}(15),\;  2\times \ce{^{1}D(5)} \\ 
     \hline
     M_{L} = 1 & 
     \begin{aligned}
     1 &= 1 + 1 - 1 + 0\\
       &= 2 + 0 + 0 - 1\\
       &= 2 + 1 - 1 - 1\\
       &= 2 + 1 + 0 - 2\\
       &= 2 + 2 - 1 - 2
   \end{aligned} & 
   \begin{aligned}
      \mathcal{D} &= 1\\
      \mathcal{D} &= 1\\
      \mathcal{D} &= 1\\
      \mathcal{D} &= 0\\
      \mathcal{D} &= 1 
   \end{aligned} & 
    \begin{cases}
    1\times (M_{S}=2)\\  
    8 \times (M_{S} = 1)\\
    14 \times (M_{S} = 0) 
    \end{cases} \implies  2\times\ce{^{3}P}(9),\ \\
 \hline 
   M_{L} = 0 & 
     \begin{aligned}
       0 &= 1 + 1 - 1 - 1\\
         &= 1 + 1 + 0 - 2\\
         &= 1 + 0 + 0 - 1\\
         &= 2 + 0 - 1 - 1\\
         &= 2 + 0 + 0 - 2\\
         &= 2 + 1 - 1 - 2\\
         &= 2 + 2 - 2 - 2 
     \end{aligned} 
     & \begin{aligned}
       \mathcal{D} &= 2\\ 
       \mathcal{D} &= 1\\
       \mathcal{D} &= 1\\
       \mathcal{D} &= 1\\
       \mathcal{D} &= 1\\
       \mathcal{D} &= 0\\
       \mathcal{D} &= 2
     \end{aligned} & \begin{dcases}
      1 \times (M_{S} = 2) \\
      8\times (M_{S} = 1)\\
      16\times(M_{S}=0) 
   \end{dcases}  \implies 2\times\ce{^{1}S(1)} \\ 
      \end{array}
\end{align*}
\begin{align*}
 g_{T} = 13 + 33 + 29 + 2 \times 9 + 2 \times 21 + 7 + 25 + 15 + 5 + 2 \times 9 + 2 = 210 = \binom{10}{4}
\end{align*}
Por tanto, los términos L-S son: 
\begin{align*}
  \boxed{\text{ T.E.: } \begin{aligned}[t]
    &\ce{^{1}I}(13),\;  \ce{^{3}H}(33),\;  \ce{^{3}G}(27),\;  2\times\ce{^{1}G}(9),\;  2\times\ce{^{3}F}(21),\;  \ce{^{1}F}(7),\;  \ce{^{5}D}(25),\;  \ce{^{3}D}(15),\\ 
    &\quad 2\times\ce{^{1}D}(5),\;  2\times\ce{^{3}P}(9),\;  2\times\ce{^{1}S}(1)
  \end{aligned}}
\end{align*}
\textbf{Nota:} Operativamente, es conveniente asignar un vector a cada regla de la siguiente forma: 
\begin{align*}
  \mathcal{D} \to  (\mathrm{mult.}( M_{S} = 2),\; \mathrm{mult.} (M_{S} = 1),\; \mathrm{mult.} (M_{S} = 0)) \implies 
  \begin{dcases}
    \mathcal{D} = 2 \to  (0,0,1)\\ 
    \mathcal{D} = 1 \to  (0,1,2)\\ 
    \mathcal{D} = 0 \to  (1,4,6)
  \end{dcases}
\end{align*}
De esta forma, para un cierto $M_{L}$ se le asigna un vector con el que podemos operar para obtener los términos espectrales. Por ejemplo, para $M_{L} = 2$ tendremos el vector asociado:
\begin{align*}
2 \times (0,0,1) + 2 \times (0,1,2) + 1 \times (1,4,6) = (1,6,12)
\end{align*}
donde los coeficientes de los vectores vienen dados por el número de veces que una regla se cumple. Así, podemos restar a este vector el asociado al anterior $M_{L}$, es decir, $M_{L} = 3$ y obtenemos: 
\begin{align*}
(1,6,12) - (0,4,8) = (1,2,4)
\end{align*}
Y este vector es lo que necesitamos para obtener los términos espectrales emergentes. Con la primera componente podemos asegurar que emerge un término con $S = 2$. De la segunda componente (2), al restar su anterior componente (1), nos da un término emergente con $S = 1$. Y con la tercera componente (4), al restar su anterior componente (2) nos da 2 términos emergentes con $S = 0$. Este mecanismo de resta incorpora las consideraciones de que términos con $S$ superior también han de tenerse en cuenta en la contabilidad de $M_{S}$ inferiores.   
 
\section{Términos j-j}\label{sec: terminos j-j} 
% !TEX root = main.tex
El acoplamiento j-j se centra en acoplar los momentos angulares individuales $j_{i}$ de cada electrón, en lugar de los momentos orbitales $l_{i}$ y spins $s_{i}$ por separado como los términos L-S. De cara a las ventajas de procedimiento que presentamos en este trabajo, seguiremos teniendo particiones, pero estas serán de $M_{J}$. No obstante, las reglas dadas por el principio de exclusión sobre las particiones serán, en comparación, más estrictas. Pues en este caso: 
\begin{center}
\fbox{
\begin{minipage}[t]{0.8\textwidth}
  \centering
 Dada una partición de $M_{J}$, entonces, en ella no podemos ningún sumando 
\end{minipage} 
} 
\end{center} 
Un apunte a tener en cuenta es que, partiendo de una cierta configuración de electrones con el mismo momento angular orbital $l$, en el momento en el que acoplamos el spin, estos electrones pueden tener $j = l + \frac{1}{2}$ o $j = l - \frac{1}{2}$. Por tanto, solo para las combinaciones donde los electrones tengan el mismo $j$, se ha de considerar el principio de exclusión. De lo contrario, los electrones son inequivalentes. Por tanto, lo primero es plantear las posibles combinaciones de $j$ de los electrones y, para aquella que sea de nuestro interés, aplicar el procedimiento de acoplamiento. \par 
Veamos un ejemplo general. Partamos de una configuración $l^{N}$ con $N < 1/2 (2l+1)\times(2s+1)$. Acoplando spin y momento angular orbital de cada electrón, los valores que toma $j$ son $j \in \qty{j_{-}, j_{+}} = \qty{l - 1/2, l + 1/2}$ y, por tanto, las posibles combinaciones de $j$  para los $N$ electrones son: 
\begin{align*}
\qty(j_{1}, j_{2}, \dots, j_{N}) \in \left\{\qty(j_{-}, j_{-}, \dots, j_{-}), \qty(j_{-}, j_{-}, \dots, j_{+}), \dots\right.\\ 
\left. \dots\qty(j_{-}, j_{+}, \dots,j_{+}), (j_{+}, j_{+}, \dots, j_{+})\right\}
\end{align*}
Siendo directo percatarse de que la cardinalidad del conjunto es de $N+1$. Para cada una de estas combinaciones podemos estudiar sus degeneraciones. Por ejemplo, para el primer y el último elemento de la tupla la degeneración estará dada por: 
\begin{align*}
g(j, j, \dots, j) = \binom{2j+1}{N} 
\end{align*}
En cambio, para una tupla con $k$ elementos $j_{+}$ y $N-k$ elementos $j_{-}$, la degeneración será:
\begin{align*}
g(\underbrace{j_{-}, j_{-}, \dots, j_{-}}_{N-k}, \underbrace{j_{+}, j_{+}, \dots, j_{+}}_{k}) = \binom{2j_{-}+1}{N-k} \times \binom{2j_{+}+1}{k}
\end{align*}
Y entonces la degeneración acumulada de \eqref{eq: numero de combinaciones final} se verifica (véase \href{https://en.wikipedia.org/wiki/Vandermonde%27s_identity}{identidad de Vandermonde}): 
\begin{align*}
\sum_{k=0}^{N} g(\underbrace{j_{-}, j_{-}, \dots, j_{-}}_{N-k}, \underbrace{j_{+}, j_{+}, \dots, j_{+}}_{k}) = \sum_{k=0}^{N} \binom{2j_{-}+1}{N-k} \times \binom{2j_{+}+1}{k} = \binom{2(2l+1)}{N}
\end{align*}
Con estas consideraciones, el procedimiento para obtener los términos j-j es análogo que para el caso Russell--Saunders. Veámoslo con los siguientes ejemplos:
 
\subsection{Ejemplo: 3 \texorpdfstring{$e^{-}$}{e^-}}
% !TEX root = main.tex
Dada la configuración $f^{3}$ estudiemos los términos j-j que podemos obtener. En primer lugar, $j \in \qty{5/2, 7/2}$ y por tanto, las combinaciones de $j$ para los 3 electrones son el conjunto: 
\begin{align*}
\qty(j_{1}, j_{2}, j_{3}) \in \left\{\qty(5/2, 5/2, 5/2),\;  \qty(5/2, 5/2, 7/2),\;  \qty(5/2, 7/2, 7/2),\;  \qty(7/2, 7/2, 7/2)\right\}
\end{align*}
La degeneración de cada combinación es: 
\begin{align*}
  g\qty(5/2, 5/2, 5/2) &= \binom{6}{3} = 20 &
g\qty(5/2, 5/2, 7/2) &= \binom{6}{2} \times  \binom{8}{1} = 120 \\
g\qty(5/2, 7/2, 7/2) &= \binom{6}{1} \times  \binom{8}{2} = 168 & 
g\qty(7/2, 7/2, 7/2) &= \binom{8}{3} = 56
\end{align*}
Y, la degeneración total es: 364. Resultado equivalente a $\binom{14}{3}$ como era de esperar. \par 
En este tipo de ejercicios, se suele preguntar por los términos j-j de la combinación de mayor degeneración. En este caso es la combinación $\qty(5/2, 7/2, 7/2)$. En esta combinación tenemos 2 electrones equivalentes con $j = 7/2$ y un electrón inequivalente con $j = 5/2$. Para tratarlo, primero acoplamos los electrones equivalentes, llamando $M_{J23}$ a la tercera componente de estos electrones acoplados. 
\begin{align*}
 \begin{array}{c | c | c}
   2M_{J23} & \text{ Particiones } & \text{ Términos }\\
   \hline 
   2M_{J23} = 12 & 12 = 7 + 5  & (7/2, 7/2)_{6} \\ 
   \hline 
    2M_{J23} = 10 & 10 = 7 + 3 & \text{ No hay nuevo término }\\ 
    \hline 
    2M_{J23} = 8 & \begin{aligned}
     8 &= 7 + 1\\
       &= 5 + 3
   \end{aligned} & (7/2, 7/2)_{4} \\
   \hline 
    2M_{J23} = 6 & \begin{aligned}
      6 &= 7 - 1\\
        &= 5 + 1\\
        \end{aligned} & \text{ No hay nuevo término }\\ 
        \hline
    2M_{J23} = 4 & \begin{aligned}
      4 &= 7 - 3\\
        &= 5 - 1\\
        &= 3 + 1
        \end{aligned}
     & (7/2, 7/2)_{2} \\
    \hline 
    2M_{J23} = 2 & \begin{aligned}
      2 &= 7 - 5\\ 
        &= 5 - 3\\
        &= 3 - 1\\ 
        \end{aligned} & \text{ No hay nuevo término }\\
        \hline 
    2M_{J23} = 0 & \begin{aligned}
      0 &= 7 - 7\\
        &= 5 - 5\\
        &= 3 - 3\\
        &= 1 - 1
    \end{aligned} & (7/2, 7/2)_{0}
\end{array}
    \end{align*}
    Donde recordemos la notación: $(j_{2},j_{3})_{J_{23}}$ . A continuación, acoplamos el electrón inequivalente con $j = 5/2$. Para ello simplemente descoponemos: 
    \begin{align*}
      J &= J_{23} \otimes j_{1} \implies \begin{cases}
        J = 6 \otimes 7/2 = 19/2 \oplus 17/2 \oplus 15/2 \oplus 13/2 \oplus 11/2 \oplus 9/2\\
      J = 4 \otimes 7/2 = 15/2 \oplus 13/2 \oplus 11/2 \oplus 9/2 \oplus 7/2 \oplus 5/2\\
      J = 2 \otimes 7/2 = 9/2 \oplus 7/2 \oplus 5/2 \oplus 3/2 \oplus 1/2 \\ 
      J = 0 \otimes 7/2 = 7/2
      \end{cases} 
    \end{align*}
    Por tanto, en los términos j-j de esta combinación son: 
    \begin{align*}
      \left\{ &(5/2, 7/2, 7/2)_{19/2}, \; (5/2, 7/2, 7/2)_{17/2}, \; 2 \times (5/2, 7/2, 7/2)_{15/2}, \; 2 \times (5/2, 7/2, 7/2)_{13/2},\right. \\ 
      &\; 2 \times (5/2, 7/2, 7/2)_{11/2}, \; 3 \times (5/2, 7/2, 7/2)_{9/2}, \; 3 \times (5/2, 7/2, 7/2)_{7/2}, \; 2 \times (5/2, 7/2, 7/2)_{5/2},\\ 
      &\left. \; (5/2, 7/2, 7/2)_{3/2}, \; (5/2, 7/2, 7/2)_{1/2} \right\}
    \end{align*}
    O en una notación más compacta: 
    \begin{align*}
      \boxed{
      (5/2, 7/2, 7/2)_\qty{19/2,\; 17/2,\;  2 \times 15/2,\;  2 \times 13/2,\;  2 \times 11/2, \; 3 \times 9/2,\;  3 \times 7/2,\;  2 \times 5/2,\;  3/2,\;  1/2}}
    \end{align*}
    \textbf{Nota:} En este caso, para calcular los términos j-j de los 2 $e^{-}$ es más sencillo estudiar la paridad de intercambio de los coeficientes de Clebsch-Gordan que acoplan $j_{2}$ y $j_{3}$ a $J_{23}$. Esta paridad viene dada por: 
    \begin{align*}
      P_{23} = (-1)^{J_{23} - j_{2} - j_{3}} = (-1)^{J_{23} - 1}
    \end{align*}
  Como $0 \leq J_{23} \leq 7$ entonces, simplemente podemos evaluar la anterior paridad para cada $J_{23}$ posible y descartar los términos pares que no cumplen la condición de antisimetría de los fermiones. Es directo darse cuenta que dichos términos son aquellos con $J_{23}$ par. Este método es más directo y nos permite comprobar nuestro resultado, pero es solo aplicable para el caso de 2 $e^{-}$ equivalentes. \par 
  La degeneración acumulada de este término es: 
\begin{align*}
  g_{T} &= 20 + 18 + 16 + 2 \times 14 + 2 \times 12 + 3 \times 10 + 3 \times 8 + 2 \times 6 + 4 = 168
\end{align*}    
Que es el resultado esperado.\par 
Por otro lado, por completitud estudiemos el caso de la combinación  hola $(7/2, 7/2, 7/2)$ donde tenemos 3 electrones equivalentes. Procedemos de la misma forma (véase el cálculo adjunto) y validamos el resultado con la degeneración acumulada: 
    \begin{align*}
    g_{T} = 16 + 12 + 10 + 8 + 6 + 4 = 56
    \end{align*}
 Que corresponde con la esperada. Así que finalmente los términos j-j de esta combinación: 
 \begin{align*}
   \boxed{
   (7/2, 7/2, 7/2)_\qty{15/2,\; 11/2,\;  9/2,\;  7/2,\;  5/2, \; 3/2,\;  3 \times 7/2,\; 5/2,\;  3/2}}
 \end{align*}

\begin{align*}
 \begin{array}{c | c | c}
   2M_{J} & \text{ Particiones } & \text{ Términos }\\
   \hline 
   2M_{J} = 15 & 15 = 7 + 5 + 3  & (7/2, 7/2, 7/2)_{15/2} \\ 
   \hline 
    2M_{J} = 13 & 13 = 7 + 5 + 1 & \text{ No hay nuevo término }\\ 
    \hline 
    2M_{J} = 11 & \begin{aligned}
     11 &= 7 + 5 -1\\
       &= 7 + 3 + 1
   \end{aligned} & (7/2, 7/2, 7/2)_{11/2} \\
   \hline 
    2M_{J} = 9 & \begin{aligned}
      9 &= 7 + 5 -3\\
        &= 7 + 3 -1\\
        &= 5 + 3 + 1\\
        \end{aligned} & (7/2, 7/2, 7/2)_{9/2}\\ 
        \hline
    2M_{J} = 7 & \begin{aligned}
      7 &= 7 + 3 - 3\\
        &= 7 + 5 - 5\\
        &= 7 + 1 - 1\\
        &= 5 + 3 - 1\\
        \end{aligned}
     & (7/2, 7/2, 7/2)_{7/2} \\
    \hline 
    2M_{J} = 5 & \begin{aligned}
      5 &= 5 + 3 - 3\\ 
        &= 7 + 3 - 5\\ 
        &= 7 + 1 - 3\\ 
        &= 7 + 5 - 7\\ 
        &= 5 + 1 - 1\\ 
        \end{aligned} & (7/2,7/2,7/2)_{5/2}\\
        \hline 
    2M_{J} = 3 & \begin{aligned}
      3 &= 5 + 1 - 3\\ 
        &= 5 + 3 - 5\\ 
        &= 7 + 1 - 5\\ 
        &= 7 + 3 - 7\\ 
        &= 7 - 1 - 3\\ 
        &= 3 + 1 - 1\\ 
        \end{aligned} & (7/2,7/2,7/2)_{3/2}\\
        \hline 
    2M_{J} = 1 & \begin{aligned}
      1 &= 3 + 1 - 3\\ 
        &= 5 - 1 - 3\\ 
        &= 5 + 1 - 5\\ 
        &= 5 + 3 - 7\\ 
        &= 7 - 1 - 5\\
        &= 7 + 1 - 7\\
        \end{aligned} & \text{ No hay nuevo término }\\
        \hline 
\end{array}
    \end{align*}
\begin{minipage}{0.37\textwidth} 
 \end{minipage} 

\subsection{Ejemplo: 4 \texorpdfstring{$e^{-}$}{e^-}}
\end{document} 

